% Unit 0: What are Differential Equations
% Practice Problems and Solutions (Part II and Part III only)
% Prepared by Staples Education
% Content informed by the local curriculum modules (0.1 to 0.4) and the
% QUIZ_DATA micro-practice and unit-mastery items for Unit 0.

\documentclass[11pt]{article}

\usepackage[margin=1in]{geometry}
\usepackage{amsmath}
\usepackage{amssymb}
\usepackage{xcolor}
\usepackage[most]{tcolorbox}
\usepackage{enumitem}
\usepackage{hyperref}
\hypersetup{colorlinks=true, linkcolor=headerblue, urlcolor=headerblue}

\tcbuselibrary{skins}

% ---------------------------------------------------------------------------
% Unified style-guide palette (see LATEX_STYLE_GUIDE.md)
% ---------------------------------------------------------------------------
\definecolor{headerblue}{HTML}{1C569E}   % overview / structural framing
\definecolor{accentgreen}{HTML}{2E7D32}  % algorithms / methods
\definecolor{accentorange}{HTML}{E27A1E} % formulas / concept takeaways
\definecolor{workpurple}{HTML}{A020F0}   % worked-example tracking (vibrant)
\definecolor{warnred}{HTML}{C0392B}      % crimson pitfall / warning

% ---------------------------------------------------------------------------
% Subsection typography: headerblue numbered sub-module titles
% ---------------------------------------------------------------------------
\usepackage{titlesec}
\titleformat{\subsection}
  {\normalfont\large\bfseries\color{headerblue}}{\thesubsection}{1em}{}

% ---------------------------------------------------------------------------
% Six core custom boxes (unbreakable, title={#1} protected)
% ---------------------------------------------------------------------------

% 1. Blue overview capsule for the topics summary.
\newtcolorbox{topicsbox}{
    enhanced,
    colback=headerblue!6!white, colframe=headerblue,
    coltitle=white, fonttitle=\bfseries,
    title=Unit 0 at a Glance: Core Sub-Topics,
    boxrule=0.6pt, arc=2mm,
    left=3mm, right=3mm, top=2mm, bottom=2mm
}

% 2. Orange formula container for standard forms and master formulas.
\newtcolorbox{formulabox}[1][Master Formula]{
    enhanced,
    colback=accentorange!8!white, colframe=accentorange,
    coltitle=white, fonttitle=\bfseries,
    title={#1},
    boxrule=0.6pt, arc=1mm,
    left=3mm, right=3mm, top=2mm, bottom=2mm
}

% 3. Green algorithm block for explicitly laid-out procedures.
\newtcolorbox{methodbox}[1][Method]{
    enhanced,
    colback=accentgreen!6!white, colframe=accentgreen,
    coltitle=white, fonttitle=\bfseries,
    title={#1},
    boxrule=0.6pt, arc=1mm,
    left=3mm, right=3mm, top=2mm, bottom=2mm
}

% 4. Purple west-bordered worked example for step-by-step solutions.
\newtcolorbox{examplebox}[1][Worked Example]{
    enhanced,
    colback=workpurple!4!white, colframe=workpurple,
    coltitle=white, fonttitle=\bfseries,
    title={#1},
    boxrule=0.4pt, sharp corners,
    borderline west={3pt}{0pt}{workpurple},
    left=5mm, right=3mm, top=2mm, bottom=2mm
}

% 5. Orange thin-bordered final concept takeaway.
\newtcolorbox{conceptbox}[1][Key Takeaway]{
    enhanced,
    colback=accentorange!5!white, colframe=accentorange,
    coltitle=white, fonttitle=\bfseries,
    title={#1},
    boxrule=0.3pt, arc=1mm,
    left=3mm, right=3mm, top=2mm, bottom=2mm
}

% 6. Crimson west-bordered warning box for common pitfalls.
\newtcolorbox{warningbox}[1][Common Pitfall]{
    enhanced,
    colback=red!3!white, colframe=warnred,
    coltitle=white, fonttitle=\bfseries,
    title={#1},
    boxrule=0.4pt, sharp corners,
    borderline west={3pt}{0pt}{warnred},
    left=5mm, right=3mm, top=2mm, bottom=2mm
}

\setlist[enumerate]{itemsep=4pt, topsep=4pt}

% Number and index sections plus subsections (e.g., 0.1) two levels deep.
\setcounter{secnumdepth}{2}
\setcounter{tocdepth}{2}

\begin{document}

% ---------------------------------------------------------------------------
% Title block
% ---------------------------------------------------------------------------
\begin{center}
    {\LARGE \bfseries Unit 0: What are Differential Equations}\\[6pt]
    {\large Prepared by Staples Education}\\[4pt]
    {\itshape Practice Problems and Solutions}
\end{center}

\vspace{0.6em}
\hrule
\vspace{0.6em}

% ===========================================================================
% PART II: REVIEW GUIDE (PRACTICE QUESTIONS)
% No \clearpage here: a standalone Practice Set opens directly on its problems.
% ===========================================================================
\section*{Part II --- Review Guide: Practice Problems}

Work the following ten problems in order. They are graded from reading and
classifying an equation, through verifying candidate solutions and applying
initial conditions, up to two full physical models. Complete solutions follow
in Part III.

\begin{enumerate}[leftmargin=2em]
    \item Verify that $y = e^{3x}$ is a solution of $\dfrac{dy}{dx} = 3y$.

    \item State the order of each equation:
          \textbf{(a)} $\dfrac{dy}{dx} + y^{2} = x$, \quad
          \textbf{(b)} $\dfrac{d^{2}y}{dx^{2}} + 3\dfrac{dy}{dx} + 2y = 0$, \quad
          \textbf{(c)} $\dfrac{d^{3}y}{dx^{3}} + \left(\dfrac{dy}{dx}\right)^{5} = \sin x$.

    \item Classify each equation as linear or nonlinear, and name the offending
          term when it is nonlinear:
          \textbf{(a)} $\dfrac{dy}{dx} + x^{2}y = e^{x}$, \quad
          \textbf{(b)} $\dfrac{dy}{dx} = y^{2}$, \quad
          \textbf{(c)} $\dfrac{d^{2}y}{dx^{2}} + \sin y = 0$.

    \item The general solution of $\dfrac{dy}{dx} = 2x$ is $y = x^{2} + C$.
          Verify it, then find the particular solution satisfying $y(1) = 5$.

    \item Show that $y = Ce^{-2x}$ satisfies $\dfrac{dy}{dx} + 2y = 0$ for every
          constant $C$, then find the $C$ for which $y(0) = 4$.

    \item Verify that $y = \cos 2x$ is a solution of the second-order equation
          $\dfrac{d^{2}y}{dx^{2}} + 4y = 0$.

    \item Classify each as an ordinary or a partial differential equation:
          \textbf{(a)} $\dfrac{dy}{dt} = ky$, \quad
          \textbf{(b)} $\dfrac{\partial u}{\partial t} = \dfrac{\partial^{2}u}{\partial x^{2}}$, \quad
          \textbf{(c)} $\dfrac{d^{2}x}{dt^{2}} + x = 0$.

    \item The general solution of $\dfrac{dy}{dx} = x + y$ is
          $y = Ce^{x} - x - 1$. Verify it, then find the particular solution
          with $y(0) = 2$.

    \item \textbf{(Application)} A population grows at a rate proportional to its
          size, $\dfrac{dP}{dt} = kP$, with $P(0) = 500$. If the population
          doubles in $4$ hours, find $k$ and write $P(t)$ explicitly.

    \item \textbf{(Application)} A radioactive sample obeys
          $\dfrac{dN}{dt} = -\lambda N$ with $N(0) = N_{0}$. Verify that
          $N(t) = N_{0}e^{-\lambda t}$ is a solution, then find $\lambda$ given
          a half-life of $10$ years.
\end{enumerate}

\vspace{0.4em}
\hrule
\vspace{0.6em}

% ===========================================================================
% PART III: SOLUTIONS
% \clearpage so the solutions print on their own page.
% ===========================================================================
\clearpage
\section*{Part III --- Complete Solutions}

\noindent\textbf{Solution 1.}\quad Differentiate the candidate and compare both
sides.
\[ \textcolor{workpurple}{\frac{dy}{dx} = 3e^{3x}} \qquad\text{and}\qquad
   \textcolor{workpurple}{3y = 3e^{3x}}. \]
Both sides equal $\textcolor{workpurple}{3e^{3x}}$ for all $x$, so $y = e^{3x}$
satisfies the equation.

\medskip
\noindent\textbf{Solution 2.}\quad The order is the highest derivative present.
\[ \textcolor{workpurple}{\text{(a) first order}}, \quad
   \textcolor{workpurple}{\text{(b) second order}}, \quad
   \textcolor{workpurple}{\text{(c) third order}}. \]
In (c) the fifth \emph{power} on $\frac{dy}{dx}$ is a degree, not an order ---
the highest derivative is still the third.

\medskip
\noindent\textbf{Solution 3.}\quad An equation is linear when $y$ and its
derivatives appear only to the first power and are never composed inside another
function.
\[ \textcolor{workpurple}{\text{(a) linear}}, \quad
   \textcolor{workpurple}{\text{(b) nonlinear: the term } y^{2}}, \quad
   \textcolor{workpurple}{\text{(c) nonlinear: the term } \sin y}. \]

\medskip
\noindent\textbf{Solution 4.}\quad Differentiate to verify, then apply the
condition.
\[ \textcolor{workpurple}{\frac{d}{dx}\big(x^{2} + C\big) = 2x} \qquad\checkmark \]
Apply $y(1) = 5$:
\[ \textcolor{workpurple}{5 = (1)^{2} + C \;\Longrightarrow\; C = 4}. \]
The particular solution is $\textcolor{workpurple}{y = x^{2} + 4}$.

\medskip
\noindent\textbf{Solution 5.}\quad Differentiate $y = Ce^{-2x}$ and substitute.
\[ \textcolor{workpurple}{\frac{dy}{dx} = -2Ce^{-2x}} \]
\[ \textcolor{workpurple}{\frac{dy}{dx} + 2y = -2Ce^{-2x} + 2Ce^{-2x} = 0} \qquad\checkmark \]
The equation holds for every $C$. Apply $y(0) = 4$:
\[ \textcolor{workpurple}{y(0) = C = 4 \;\Longrightarrow\; y = 4e^{-2x}}. \]

\medskip
\noindent\textbf{Solution 6.}\quad Differentiate twice and substitute.
\[ \textcolor{workpurple}{\frac{dy}{dx} = -2\sin 2x}, \qquad
   \textcolor{workpurple}{\frac{d^{2}y}{dx^{2}} = -4\cos 2x} \]
\[ \textcolor{workpurple}{\frac{d^{2}y}{dx^{2}} + 4y = -4\cos 2x + 4\cos 2x = 0} \qquad\checkmark \]
So $y = \cos 2x$ is a solution of the second-order equation.

\medskip
\noindent\textbf{Solution 7.}\quad An ordinary equation has a single independent
variable; a partial equation has several.
\[ \textcolor{workpurple}{\text{(a) ODE (only } t)}, \quad
   \textcolor{workpurple}{\text{(b) PDE (}t \text{ and } x)}, \quad
   \textcolor{workpurple}{\text{(c) ODE (only } t)}. \]

\medskip
\noindent\textbf{Solution 8.}\quad Differentiate $y = Ce^{x} - x - 1$ and check
against $x + y$.
\[ \textcolor{workpurple}{\frac{dy}{dx} = Ce^{x} - 1} \]
\[ \textcolor{workpurple}{x + y = x + \big(Ce^{x} - x - 1\big) = Ce^{x} - 1} \qquad\checkmark \]
The two agree, so the family solves the equation. Apply $y(0) = 2$:
\[ \textcolor{workpurple}{2 = Ce^{0} - 0 - 1 = C - 1 \;\Longrightarrow\; C = 3} \]
\[ \textcolor{workpurple}{y = 3e^{x} - x - 1}. \]

\medskip
\noindent\textbf{Solution 9.}\quad A rate proportional to size gives exponential
growth, $\textcolor{workpurple}{P(t) = P_{0}e^{kt}}$, and here $P_{0} = 500$:
\[ \textcolor{workpurple}{P(t) = 500\,e^{kt}}. \]
Doubling in $4$ hours means $P(4) = 1000$:
\[ \textcolor{workpurple}{1000 = 500\,e^{4k} \;\Longrightarrow\; e^{4k} = 2 \;\Longrightarrow\; 4k = \ln 2} \]
\[ \textcolor{workpurple}{k = \frac{\ln 2}{4}}. \]
Hence
\[ \textcolor{workpurple}{P(t) = 500\,e^{(\ln 2 / 4)\,t} = 500\cdot 2^{\,t/4}}. \]

\medskip
\noindent\textbf{Solution 10.}\quad First verify the candidate by
differentiating:
\[ \textcolor{workpurple}{\frac{dN}{dt} = -\lambda N_{0}e^{-\lambda t} = -\lambda N} \qquad\checkmark \]
so $N(t) = N_{0}e^{-\lambda t}$ solves the decay law. A half-life of $10$ years
means $N(10) = \tfrac{1}{2}N_{0}$:
\[ \textcolor{workpurple}{\tfrac{1}{2}N_{0} = N_{0}e^{-10\lambda} \;\Longrightarrow\; e^{-10\lambda} = \tfrac{1}{2}} \]
\[ \textcolor{workpurple}{-10\lambda = -\ln 2 \;\Longrightarrow\; \lambda = \frac{\ln 2}{10} \text{ per year}}. \]

\vspace{0.8em}
\begin{conceptbox}[Unit 0 Key Takeaway]
Every problem here rests on one habit --- treat the unknown as a
\textcolor{accentorange}{function} and let the equation test it. You
\textcolor{accentorange}{verify} a candidate by substituting it and its
derivatives and checking both sides agree for all $x$; you pass from the
\textcolor{accentorange}{general solution} to a particular curve by using an
\textcolor{accentorange}{initial condition} to fix the arbitrary constant. The
same proportional rate law drives growth and decay alike.
\end{conceptbox}

\end{document}
